\documentclass[10pt,a4paper]{amsart}
\usepackage[margin=19mm]{geometry}
\usepackage{amsmath,amssymb,mathtools}
\usepackage[scr=rsfs,cal=euler]{mathalfa}
\usepackage{xfrac}
\usepackage[unicode,hidelinks,bookmarks=false]{hyperref}
\newcommand{\ie}{i.e.\ }
\newcommand{\cf}{cf.\ }
\newcommand{\resp}{respectively\ }
\newcommand{\Subsection}{\subsection}
\providecommand{\nobreakdash}{\nobreak\hbox{-}}
\setlength{\emergencystretch}{2em}
\multlinegap0pt
\usepackage{bm}
\usepackage{bbm}
\usepackage{dsfont}
\usepackage{xspace}
\usepackage{cancel}
\usepackage{mathtools}
\usepackage{amsthm}
\makeatletter
\@ifundefined{theorem}{\newtheorem{theorem}{Theorem}}{}
\@ifundefined{lemma}{\newtheorem{lemma}[theorem]{Lemma}}{}
\makeatother
\renewcommand{\epsilon}{{\varepsilon}}
\title[Rainbow cycles in triangle-free graphs]{Rainbow cycles in triangle-free graphs}
\author{Andrzej Czygrinow, Skand Parvatikar}
\date{Source: arXiv:2606.14097v1 (CC BY 4.0)}
\begin{document}
\maketitle

\noindent\emph{Source and licence.} Rainbow cycles in triangle-free graphs. Version arXiv:2606.14097v1 (CC BY 4.0), distributed under the Creative Commons Attribution 4.0 International licence, as recorded in the arXiv licence metadata for this exact version and shown on the abstract page https://arxiv.org/abs/2606.14097v1. Full rights notice: licenses/arxiv-2606.14097-CC-BY-4.0.txt.\par\smallskip

\noindent\textit{Editorial scope.} Counted unit: the Lemma whose source label is \texttt{non-ext-lem2} in the article (source0.tex lines 345--347, source label \texttt{non-ext-lem2}), delivered together with the article's own complete original proof of it (source0.tex lines 348--402, one \texttt{proof} environment of 55 lines). No mathematics is written, repaired or re-derived: statement and proof are the article's own text line for line. Required context retained uncounted: eq3 (lines 274--277); non-ext-lem1 (lines 305--308). Every definition, display and label of those lines is retained unchanged. Omitted from this package and not redistributed: 1--273, 278--304, 309--344, 403--1148. Nothing inside a retained excerpt is omitted or altered: every retained body is delivered exactly as the article's lines are written, so no omission receipt of any kind accompanies this package. Citations: the retained excerpts carry no citation command at all, so no bibliography entry of the article is transcribed and no reference is redistributed. Reference adaptation: none was needed and none was made; every reference inside the delivered unit resolves inside it, so no unresolved reference or citation is produced. Printed slips kept literal: no printed word, symbol, hypothesis or number of the counted statement or of its proof was altered. Layout and class: the article's own class is \texttt{article} with packages including inputenc, amsmath, amsthm, xcolor, amssymb, mathrsfs, tabularx, booktabs, enumitem, lineno, authblk, mathtools, graphics, graphicx; the wrapper is the lane's pinned \texttt{amsart} editorial wrapper at 10pt on A4, which restates the theorem-like environments the retained text needs and adds the article's own macro definitions that the retained text needs (\texttt{\textbackslash epsilon}), with extra wrapper packages (bm, bbm, dsfont, xspace, cancel, mathtools). The delivered numbering is the unit's own. Assets: no retained span contains a figure environment, a graphics-inclusion command, a PDF-page inclusion, a table or a TikZ picture; no image file is copied into this package, the package declares no asset dependency and no image file is redistributed with it. Licence: version arXiv:2606.14097v1 of this article is distributed under the Creative Commons Attribution 4.0 International licence, as recorded in the arXiv licence metadata for this exact version and shown on the abstract page https://arxiv.org/abs/2606.14097v1. A full rights notice is delivered at licenses/arxiv-2606.14097-CC-BY-4.0.txt. Characters: every retained excerpt is pure ASCII, so the delivered engine, XeLaTeX with its default Latin Modern text font, reports no missing character.
\begin{equation}\label{eq3}
    |N^+_G(u)|+|N^+_G(v)|\geq 2\cdot(1/5-\epsilon)m - |C|\geq
    (1/5- 6\epsilon)m.
\end{equation}

\begin{lemma}\label{non-ext-lem1}
 $||G||\geq |W_0|\cdot |W_1|-\sqrt{\xi}m^2.$
 %For all but at most $\sqrt{\xi}m^2$ pairs $(w_0, w_1)\in W_0\times W_1$, $\{w_0,w_1\}\in E(\widetilde{G}).$
\end{lemma}

\begin{lemma}\label{non-ext-lem2}
$|E_D(B_1, B_0)|\geq (1-\mu)|B_0|\cdot |B_1|$ where $\mu=O(\xi^{1/8}/\sqrt{\eta})$. 
\end{lemma}

\begin{proof}
In view of Lemma \ref{non-ext-lem1},
$$|E_G(W_1, W_0)|+|E_G(W_0,W_1)|\geq |W_0|\cdot |W_1|- \sqrt{\xi} m^2.$$
Let $\nu:=\frac{\xi^{1/4}}{\eta}$ and note that $\nu \ll \eta$. Let $W_i'= \{v\in W_i: |W_{1-i} \cap (N^+_G(v)\cup N^-_G(v))| \geq (1-\nu)|W_{1-i}|\}.$
We have
$$|W_1|\cdot |W_0| -\sqrt{\xi} m^2\leq |W_i'|\cdot |W_{1-i}|+ (1-\nu)|W_i\setminus W_i'|\cdot |W_{1-i}|\leq |W_1|\cdot |W_0|-\nu |W_i\setminus W_i'|\cdot |W_{1-i}|.$$
Therefore, \[
    \sqrt{\xi}m^2 \geq \nu|W_i \setminus W'_i|\cdot |W_{1 - i}|
\]
Since $|E_D(W_1,W_0)|=|E_D(B_1, B_0)|\geq \eta m^2$, $\min{|W_{1-i}|, |W_i|}\geq \eta m$, and so using the definition of $\nu$,
\begin{equation}\label{eq9}
|W_i\setminus W_i'|\leq \frac{\sqrt{\xi}m^2}{\nu |W_{1-i}|}\leq \xi^{1/4}m\leq \nu |W_i|.\end{equation}
Consequently,
$$|W_i'|\geq (1-\nu)|W_i|.$$
%Consequently,
%\begin{equation}\label{eq9}
%|W_i'|\geq 
%(1-\nu)|W_i|.
%\end{equation}
%As a result, by the definition of $W_i'$,
%$$|E_G(W_1', W_0')|+|E_G(W_0',W_1')|\geq (1-\nu)|W_1'|\cdot |W_0|- |W_1'|\cdot |W_0\setminus W_0'|\geq (1-2\nu)|W_1'|\cdot|W_0'|.$$

%Let $G':=G[W_0',W_1'].$ 
%Recall that by (\ref{eq3})  for every $uv\in E(G)$,
%\begin{equation}\label{eq10}|N^+_{G}(u)|+|N^+_{G}(v)|\geq (1/5-6\epsilon)m- \nu (|W_0|+|W_1|)\geq (1/5-\nu)m.
%\end{equation}
Let $t:= \min_{v\in W_0'}|N^+_{G}(v)|$ and suppose $x\in W_0'$  satisfies $t= |N^+_{G}(x)|.$ 
Then 
\begin{equation}\label{eq-ne-t}
    t\leq (1/5-\eta)m.
\end{equation}
Indeed, if $t>(1/5-\eta)m$, then using the fact that $G$ is an oriented graph, $|W_i|< m/4$ and (\ref{eq9}), we have

$$|E_D(W_1, W_0)|\leq |W_1|\cdot|W_0\setminus W_0'|+|W_1\setminus W_1'|\cdot |W_0'|+ (\eta+\epsilon)m|W_1'|< (\eta/2 +\nu)m^2,$$
which is less than $\eta m^2$ since $\nu \ll \eta$. 

By the definition of $x$, from (\ref{eq3}), for every $z\in N^+_{G}(x)\cup N^-_{G}(x)$,  $|N^+_{G}(z)|\geq (1/5-6\epsilon)m- t.$ 
Consequently, since $x\in W_0'$, $|N^+_{G}(x)\cup N^-_{G}(x)|\geq (1-\nu)|W_1|$, and so
\begin{equation}\label{eq-E_G}
    |E_{G}(W_1, W_0)|\geq (m/5-t -6\epsilon m)(1-\nu)|W_1|\geq (m/5- t-2\nu m)|W_1|.
\end{equation}
We now consider two cases based on the value of $t$.

If $t\leq \sqrt{\nu}m$, then from (\ref{eq-E_G}), 
$$|E_D(B_1, B_0)|\geq (m/5- 3\sqrt{\nu} m)|W_1|\geq (1- O(\sqrt{\nu}))|B_1|\cdot |B_0|,$$
because $|B_1|=|W_1|$ and $|B_0|\leq m/5+\epsilon m$.

Now suppose $t>\sqrt{\nu}m$. We will show that this case leads to a contradiction. Since $|E_G(W_1, W_0\setminus W_0')|<\nu |W_1|\cdot |W_0|$, from (\ref{eq-E_G}), there exists $v\in W_1'$ such that $|N_G^-(v)|\geq m/5 -t -3\nu m$. Let $A_1= N^-_G(v)$, $A_2=N^+_G(v)$, and $A_3=N^+_G(B)$. Then $A_1\cap A_2=\emptyset$, $|A_2|\geq t $ and since there is no $C_4$, for every $w\in A_3$, $N^+_G(w)\cap A_1=\emptyset$. Since $|W_0\setminus A_1|<t+4\nu m$, we have
$$(t+4\nu m)|A_3|> |E_G(W_0\setminus A_1, A_3)|+ |E_G(A_3, W_0\setminus A_1)|\geq |E_G(A_2, A_3)|+t|A_3|.$$
 Thus
 $$4\nu m |A_3|> |E_G(A_2,A_3)|\geq (m/5-t-6\epsilon m)|A_2|\geq (m/5-t-6\epsilon m)t.$$
However, since $\nu \ll \eta$ and $t>\sqrt{\nu}m$, by (\ref{eq-ne-t}),
$(m/5-t-6\epsilon m)t> \sqrt{\nu} m^2/6$, contradicting $\nu \ll \eta \leq 1$.

\end{proof}

\end{document}
